\documentclass[conference]{IEEEtran}
\IEEEoverridecommandlockouts

\usepackage{cite}
\usepackage{amsmath,amssymb,amsfonts}
\usepackage{algorithmic}
\usepackage{graphicx}
\usepackage{textcomp}
\usepackage{xcolor}
\usepackage{booktabs}
\usepackage{tabularx}
\usepackage{array}

\def\BibTeX{{\rm B\kern-.05em{\sc i\kern-.025em b}\kern-.08em
    T\kern-.1667em\lower.7ex\hbox{E}\kern-.125emX}}

\begin{document}

\title{Intelligent Invoice Processing for Automated Accounting Book Generation with Multimodal Table Correction and Invariant Enforcement}

\author{\IEEEauthorblockN{Anand Misra}
\IEEEauthorblockA{\textit{Dept. of Computer Science and Eng.} \\
\textit{Bharati Vidyapeeth's College of Eng.}\\
New Delhi, India \\
04011502723.anand@ipu.ac.in}
\and
\IEEEauthorblockN{Pratham Nahata}
\IEEEauthorblockA{\textit{Dept. of Computer Science and Eng.} \\
\textit{Bharati Vidyapeeth's College of Eng.}\\
New Delhi, India \\
05611502723.pratham@ipu.ac.in}
\and
\IEEEauthorblockN{Dr. Jolly Parikh}
\IEEEauthorblockA{\textit{Associate Professor, Dept. of CSE} \\
\textit{Bharati Vidyapeeth's College of Eng.}\\
New Delhi, India \\
jolly.parikh@bvcoedelhi.ac.in}
}

\maketitle

\begin{abstract}
The automated processing of business financial documents represents a fundamental challenge in computational document intelligence and enterprise automation. While modern Vision-Language Models (VLMs) have achieved high milestones in extracting text from document images, they stop at isolated field extraction. The essential cognitive leap---interpreting raw tabular data as legal financial transactions, classifying capital versus operational expenditures under governing accounting standards (AS 10 / Ind AS 16), resolving statutory Indian GST tax treatments (CGST/SGST vs IGST), and posting mathematically balanced double-entry journal entries---remains manual, fragmented, and error-prone. This paper presents an end-to-end Multimodal Document Intelligence and Double-Entry Accounting Framework that bridges this semantic chasm. The system couples template-free vision heuristic extraction with multimodal table layout correction to decouple product line items from footer summaries. An Accounting Interpretation Engine maps items to a formal Chart of Accounts, enforces AS 10 capitalization policies ($>$INR 10,000 threshold), determines Place of Supply (POS) tax rules, and evaluates Section 17(5) blocked Input Tax Credits. Before ledger commitment, a zero-tolerance Deterministic Invariant Validation Suite mathematically enforces exact decimal double-entry balance (sum of debits equals sum of credits). Evaluated across diverse GST tax invoices, the framework achieves 98.4\% Table TEDS accuracy, eliminates floating-point drift with 100\% exact decimal balance enforcement, and successfully passes 48 out of 48 comprehensive automated system tests, compiling real-time Trial Balance and Balance Sheet statements.
\end{abstract}

\begin{IEEEkeywords}
Document AI, Intelligent Invoice Processing, Table Layout Correction, Double-Entry Bookkeeping, Invariant Validation, Indian GST, Capitalization Threshold, Automated Accounting Books.
\end{IEEEkeywords}

\section{Introduction}
The rapid digitization of global business operations has transformed how organizations exchange financial documents. Invoices, purchase orders, and debit notes constitute the fundamental transactional heartbeat of commerce. In India, the implementation of the Goods and Services Tax (GST) framework introduced rigorous digital compliance requirements, necessitating meticulous tracking of B2B transactions, HSN/SAC classifications, Input Tax Credit (ITC) eligibility, and state-level tax segregations (CGST, SGST, and IGST). Despite these technological advancements in document transmission, the downstream accounting workflow---ingesting invoices, interpreting economic reality, applying governing accounting standards, and posting verified double-entry journal entries---remains overwhelmingly manual, fragmented, and vulnerable to costly human error \cite{rose2020zero}.

\subsection{The Industrial Problem: Manual Accounting Bottlenecks}
In micro, small, and medium enterprises (MSMEs) and corporate accounting practices, finance professionals spend countless hours manually reading incoming supplier bills, verifying line-item computations, matching purchase orders, checking tax calculations, and typing entries into enterprise resource planning (ERP) software such as Tally or Zoho Books. This manual paradigm suffers from severe limitations: high labor latency (averaging 15 to 25 minutes per complex invoice), vulnerability to transposed digits and calculation oversights, ambiguity in choosing Chart of Accounts (COA) codes, and complex compliance overhead in verifying 15-character Goods and Services Tax Identification Numbers (GSTIN), Place of Supply (POS) state codes, and Section 17(5) statutory blocked input tax credits \cite{rana2025template}.

\subsection{The Semantic and Accounting Chasm in Document AI}
Over the past half-decade, Document AI has witnessed tremendous breakthroughs. Multimodal transformer models such as LayoutLMv1-v3, Donut, and modern Vision-Language Models (VLMs) have demonstrated exceptional capabilities in extracting bounding boxes, text tokens, and key-value pairs from scanned invoices \cite{yao2026multimodal, xu2020layoutlm, xu2021layoutlmv2, huang2022layoutlmv3}. However, a profound disconnect exists between document extraction and financial bookkeeping. Existing Document AI systems operate under an ``extraction-only'' philosophy: they return structured JSON payloads containing raw strings and numerical values, leaving the entire cognitive burden of accounting interpretation to human accountants.

Specifically, current models fail to bridge the semantic and accounting chasm in four critical dimensions:
\begin{enumerate}
    \item \textit{Economic Reality Interpretation:} An invoice line reading ``Dell XPS 15 Laptop - INR 1,15,000'' is not merely an expense; it represents a capital asset acquisition that must be capitalized to the balance sheet rather than expensed to the profit and loss statement.
    \item \textit{Capitalization Threshold Enforcement:} Accounting standards (such as AS 10 / Ind AS 16) mandate that items exceeding an organizational threshold (e.g., INR 10,000) be recognized as Property, Plant, and Equipment (PPE), whereas lower-value durable items (e.g., a INR 1,250 mouse) must be expensed immediately.
    \item \textit{Statutory Tax Disambiguation:} Under Indian GST law, taxes must be bifurcated into Central GST (CGST) and State GST (SGST) for intra-state transactions, or Integrated GST (IGST) for inter-state supplies, while blocked credits under Section 17(5) (e.g., corporate catering) must be expensed rather than claimed as asset credits.
    \item \textit{Invariant Double-Entry Balance:} Raw extracted text frequently contains OCR rounding errors. Posting unbalanced entries directly to an accounting ledger violates the fundamental equation of accounting, corrupting financial integrity \cite{pacioli1494summa}.
\end{enumerate}

\subsection{The Double-Entry Discipline and Mathematical Invariants}
In double-entry bookkeeping, established by Luca Pacioli in 1494, every financial event requires dual entries: for every debit, an exact matching credit must exist. In formal vector notation, let $\mathcal{D}$ and $\mathcal{C}$ represent the sets of debit and credit ledger lines for a transaction. The fundamental invariant states that the sum of debits must strictly equal the sum of credits with zero numerical tolerance:
\begin{equation}
\Delta_{\text{bal}} = \left| \sum_{j \in \mathcal{D}} D_j - \sum_{k \in \mathcal{C}} C_k \right| = 0.000000
\end{equation}
Standard machine learning extraction pipelines treat numerical fields as independent scalar predictions, frequently violating this law. A production accounting system requires deterministic invariant guards that mathematically certify transaction balance prior to database commit.

\subsection{Primary Contributions}
To resolve these challenges, this paper presents an end-to-end Multimodal Document Intelligence and Double-Entry Accounting Framework. The primary technical contributions are:
\begin{enumerate}
    \item A template-free multimodal extraction and table layout correction pipeline (adapting Yao et al., 2026 \cite{yao2026multimodal}) that decouples product line items from trapped footer summaries and resolves multi-column alignment errors.
    \item An Accounting Interpretation Engine that classifies line items to a formal Chart of Accounts, enforces AS 10 capitalization policies ($>$INR 10,000 threshold), determines statutory GST Place of Supply rules, and identifies Section 17(5) blocked Input Tax Credits.
    \item A zero-tolerance Deterministic Invariant Validation Suite that mathematically enforces unit math, tax consistency, subtotal reconciliation, and exact decimal double-entry balance with zero floating-point drift.
    \item An authoritative general ledger posting engine and real-time financial statement compiler (Trial Balance, P\&L, Balance Sheet) evaluated on 35 automated test cases with a 100\% pass rate.
\end{enumerate}

\section{Related Work}
\subsection{Evolution of Document Intelligence}
Automated document understanding has evolved through three historical paradigms: traditional optical character recognition (OCR) paired with rigid heuristic regex rules; spatial-aware multimodal transformers (LayoutLMv1-v3 \cite{xu2020layoutlm, xu2021layoutlmv2, huang2022layoutlmv3}); and OCR-free vision transformers (Donut \cite{kim2022donut}, Nougat \cite{blecher2023nougat}). While multimodal transformers capture 2D spatial layouts and visual cues, they remain confined to token-level classification. They lack domain-specific accounting reasoning and fail to understand how extracted figures interact within legal ledgers.

\subsection{Table Structure Recognition \& Multimodal Correction}
Table extraction in business invoices is notoriously challenging due to borderless tables, multi-line product descriptions, merged columns, and subtotal summaries embedded within the tabular grid. PubTables-1M \cite{smock2022pubtables} and Table-Transformer introduced deep object detection for table boundaries, yet often fail on invoice footers where discount rows and tax summaries mimic item rows. Recently, Yao et al. (2026) proposed a multimodal table layout correction pipeline for unstructured documents, demonstrating that spatial coordinate alignment combined with layout syntax correction dramatically reduces tabular extraction errors \cite{yao2026multimodal}. Our work adapts and extends this approach to financial invoices, specifically isolating line items from taxable totals.

\subsection{Accounting Information Systems and Commercial AP Tools}
Traditional enterprise resource planning (ERP) platforms rely entirely on human data entry or rigid electronic data interchange (EDI) formats. Rule-based Accounts Payable (AP) automation tools utilize template matching, which breaks whenever a supplier alters invoice geometry. Furthermore, existing AI-based AP tools provide predictive GL coding but lack deterministic mathematical invariant guards, frequently requiring human accountants to resolve decimal rounding discrepancies \cite{vasoya2022ai}. Table \ref{tab:comparison} compares existing document processing frameworks against our proposed framework across critical functional dimensions.

\begin{table}[htbp]
\caption{Comparative Analysis of Document AI Frameworks}
\begin{center}
\begin{tabularx}{\columnwidth}{lXX}
\toprule
\textbf{Framework / Tool} & \textbf{Template-Free?} & \textbf{Double-Entry Guard?} \\
\midrule
Traditional OCR + Regex & No (Rigid Templates) & None (Manual) \\
LayoutLMv1-v3 & Partial (Token VDU) & None (Extraction Only) \\
Donut / Nougat & Yes (Vision-to-Text) & None (Field Only) \\
Commercial AP (Dext / Vic.ai) & Yes (Cloud Heuristic) & Partial (Soft Warning) \\
\textbf{Proposed Framework} & \textbf{Yes (Multimodal Corr.)} & \textbf{Strict (Zero Drift Invariants)} \\
\bottomrule
\end{tabularx}
\label{tab:comparison}
\end{center}
\end{table}

\section{System Architecture}
\subsection{End-to-End Pipeline Overview}
The proposed architecture is structured into seven sequential stages, as illustrated in Fig. \ref{fig:pipeline}. The pipeline transforms raw invoice files (vector PDFs or raster scans) into legally compliant, mathematically balanced accounting books without requiring predefined templates:
\begin{enumerate}
    \item \textit{Stage 1: Document Ingestion \& Spatial Preprocessing}
    \item \textit{Stage 2: Spatial Layout \& Vector Text Normalization}
    \item \textit{Stage 3: Table Layout Correction (Yao et al., 2026 adaptation)}
    \item \textit{Stage 4: Semantic Decoupling \& Normalized Data Extraction}
    \item \textit{Stage 5: Accounting Interpretation Engine (AS 10 \& GST POS)}
    \item \textit{Stage 6: Deterministic Invariant Validation Suite}
    \item \textit{Stage 7: General Ledger Posting \& Real-Time Financial Statements}
\end{enumerate}

\begin{figure}[htbp]
\centerline{\includegraphics[width=0.92\columnwidth]{../../report_figures/fig1_pipeline_architecture.png}}
\caption{End-to-end system architecture: Multimodal document intelligence to verified financial books.}
\label{fig:pipeline}
\end{figure}

\subsection{Multimodal Ingestion and Spatial Normalization}
Stage 1 ingests documents in PDF, PNG, or JPEG format via a dual-path preprocessor. For vector PDFs, \texttt{pymupdf4llm} extracts text tokens, normalized bounding boxes $(x_0, y_0, x_1, y_1)$, and markdown tables. For raster scans, OpenCV applies Gaussian thresholding, deskewing via Hough transform, and contour boundary detection \cite{bradski2008opencv}. All coordinates are normalized to a canonical $1000 \times 1000$ grid, ensuring scale invariance across varied document resolutions.

\section{Proposed Methodology}
\subsection{Table Layout Correction and Semantic Decoupling}
Business invoices frequently present tabular layouts where column dividers are missing, line descriptions span multiple lines, and summary rows (Subtotal, Tax, Round Off, Grand Total) are embedded within the item grid. Adapting Yao et al. (2026) \cite{yao2026multimodal}, our Table Layout Corrector applies spatial clustering and semantic keyword decoupling (Fig. \ref{fig:table_corr}).

\begin{figure}[htbp]
\centerline{\includegraphics[width=0.90\columnwidth]{../../report_figures/fig2_table_correction_process.png}}
\caption{Multimodal table layout correction and footer summary trapping mechanism (Yao et al. adaptation).}
\label{fig:table_corr}
\end{figure}

The algorithm operates in three phases:
\begin{enumerate}
    \item \textit{Column Boundary Resolution:} Horizontal projection profiles identify vertical text clusters, establishing bounds for Description, HSN/SAC, Quantity, Unit Price, and Total.
    \item \textit{Multi-line Concatenation:} Consecutive text spans lacking numeric values in quantity/price columns are merged into the preceding item description.
    \item \textit{Footer Summary Trapping:} Rows containing summary tokens (``Total'', ``CGST'', ``SGST'', ``IGST'', ``Round Off'') are severed from the line item array and populated into a \texttt{TrappedTotalsMap}, ensuring summaries are never treated as purchasable goods.
\end{enumerate}

\begin{figure}[htbp]
\centerline{\includegraphics[width=0.90\columnwidth]{../../report_figures/fig3_accounting_decision_tree.png}}
\caption{Accounting interpretation, capitalization threshold (AS 10), and statutory GST decision flowchart.}
\label{fig:decision_tree}
\end{figure}

\subsection{Accounting Interpretation Engine}
Stage 5 translates raw commercial items into legal debits and credits by operationalizing statutory accounting rules (Fig. \ref{fig:decision_tree}).

\subsubsection{Capitalization Policy under AS 10 / Ind AS 16}
Under AS 10 (Property, Plant and Equipment) \cite{icai2018as10}, entities establish a capitalization threshold. Durable items ($>$1 year utility) exceeding the threshold $\tau_{\text{cap}} = \text{INR } 10,000.00$ are capitalized as Non-Current Assets on the Balance Sheet. Low-value consumables or items with unit price $\le \text{INR } 10,000.00$ are expensed to Profit \& Loss accounts:
\begin{equation}
\text{Classify}(L_i) = \begin{cases} 
\text{Non-Current Asset (PPE)}, & \text{if } U_i > \tau_{\text{cap}} \wedge \text{Durable}(L_i) \\ 
\text{Operating Expense}, & \text{if } U_i \le \tau_{\text{cap}} \vee \neg \text{Durable}(L_i) 
\end{cases}
\end{equation}

\subsubsection{Statutory GST Place of Supply (POS) Logic}
Tax treatment is governed by supplier location and recipient Place of Supply (POS). The system extracts the two-digit state prefix from the supplier's GSTIN and compares it with the recipient's designated POS code:
\begin{equation}
\text{TaxType} = \begin{cases} 
\text{CGST } (r_i/2) + \text{SGST } (r_i/2), & \text{if } \text{State}(\text{Supp}) = \text{State}(\text{POS}) \\ 
\text{IGST } (r_i), & \text{if } \text{State}(\text{Supp}) \neq \text{State}(\text{POS}) 
\end{cases}
\end{equation}
When $\text{State}(\text{Supp}) = \text{State}(\text{POS})$, the supply is intra-state, splitting tax equally between CGST and SGST. When $\text{State}(\text{Supp}) \neq \text{State}(\text{POS})$, it is inter-state, subject exclusively to IGST. Both are mapped to current asset accounts (Codes 1081--1083) as claimable ITC.

\subsubsection{Statutory Blocked Credits under Section 17(5)}
Under Section 17(5) of the CGST Act \cite{cgst2017act}, input tax credit is disallowed for food and beverages, catering, corporate gifts, club memberships, and passenger vehicles:
\begin{equation}
\text{Account}(\text{Tax}_i) = \begin{cases} 
\text{Ineligible GST Exp (Code 5080)}, & \text{if } L_i \in \mathcal{B}_{17(5)} \\ 
\text{Input GST Asset (Codes 1081--1083)}, & \text{if } L_i \notin \mathcal{B}_{17(5)} 
\end{cases}
\end{equation}
Flagged taxes are debited to an operating expense---``Ineligible GST Expense'' (Account Code 5080)---rather than an ITC asset.

\begin{table}[htbp]
\caption{Standard Chart of Accounts (COA) Definitions}
\begin{center}
\begin{tabularx}{\columnwidth}{llX}
\toprule
\textbf{Code} & \textbf{Account Name} & \textbf{Category / Head} \\
\midrule
1010 & Cash in Hand & Current Asset \\
1050 & Office Furniture \& Fixtures & Non-Current Asset (PPE) \\
1060 & Computer \& Hardware Assets & Non-Current Asset (PPE) \\
1081 / 1082 & Input CGST / Input SGST & Current Asset (Tax Credit) \\
1083 & Input IGST (Inter-State) & Current Asset (Tax Credit) \\
2010 & Accounts Payable (Creditors) & Current Liability \\
5010 & Office Supplies Expense & Operating Expense \\
5060 & Staff Welfare Expense & Operating Expense \\
5080 & Ineligible GST Expense & Operating Expense (Blocked) \\
\bottomrule
\end{tabularx}
\label{tab:coa}
\end{center}
\end{table}

\begin{table}[htbp]
\caption{Statutory GST Place of Supply (POS) Rules}
\begin{center}
\begin{tabularx}{\columnwidth}{llX}
\toprule
\textbf{State Relation} & \textbf{Supply Type} & \textbf{Tax Accounts Applied} \\
\midrule
Supplier == Recipient & Intra-State & CGST (1081) + SGST (1082) \\
Supplier != Recipient & Inter-State & IGST (1083) \\
Import of Goods & Inter-State & IGST (1083) + Customs \\
Section 17(5) Disallowed & Any Supply & Ineligible GST Expense (5080) \\
\bottomrule
\end{tabularx}
\label{tab:gst}
\end{center}
\end{table}

\subsection{Deterministic Invariant Validation Suite}
Stage 6 executes a six-point invariant verification suite. Operating as a deterministic mathematical guard, all six conditions must hold with exact zero tolerance:
\begin{itemize}
    \item \textbf{Invariant 1 (Line Math):} $| Q_i \cdot U_i - T_i | \le 0.01$
    \item \textbf{Invariant 2 (Tax Consistency):} $\text{Tax}_i = \text{Round}(T_i \cdot r_i, 2)$
    \item \textbf{Invariant 3 (Subtotal Reconciliation):} $\sum T_i = \text{Subtotal}$
    \item \textbf{Invariant 4 (Grand Total Completeness):} $\text{GrandTotal} = \text{Subtotal} + \sum \text{Taxes} + \Delta_{\text{round}}$
    \item \textbf{Invariant 5 (COA Validity):} All account codes exist as active accounts in master COA (Table \ref{tab:coa}).
    \item \textbf{Invariant 6 (Double-Entry Balance):} $\Delta_{\text{bal}} = | \sum \text{Debits} - \sum \text{Credits} | = 0.000000$
\end{itemize}

\section{Implementation Details}
\subsection{Technology Stack \& Decimal Precision}
The backend is built in Python 3.13 using FastAPI and SQLAlchemy 2.0 ORM with typed declarative models. Standard IEEE 754 floating-point types are excluded across financial paths; monetary amounts, tax rates, and ledger balances use Python \texttt{Decimal} objects with \texttt{ROUND\_HALF\_UP} precision. This eliminates binary floating-point drift (e.g., $0.1 + 0.2 \ne 0.3$), guaranteeing exact cent-level reconciliation across all accounts \cite{ieee754}.

\begin{figure}[htbp]
\centerline{\includegraphics[width=0.92\columnwidth]{../../report_figures/fig4_database_er_diagram.png}}
\caption{Relational database Entity-Relationship (ER) schema for the autonomous accounting platform.}
\label{fig:er_diag}
\end{figure}

\subsection{Relational Schema Design \& Audit Trail}
Figure \ref{fig:er_diag} illustrates the complete Entity-Relationship (ER) schema. Referential integrity constraints enforce multi-tenant isolation across \texttt{tenants}, \texttt{invoices}, \texttt{invoice\_line\_items}, \texttt{journal\_entries}, \texttt{journal\_lines}, and \texttt{accounts}. To satisfy corporate audit standards (Companies Act, 2013 and GST rules), every state transition is recorded in an immutable \texttt{audit\_logs} table. Committed journal entries are immutable; adjustments require explicit offsetting credit/debit memorandum entries.

\section{Experimental Evaluation}
\subsection{Benchmark Dataset \& Setup}
Evaluation was conducted on public datasets (CORD \cite{park2019cord}, SROIE \cite{huang2019sroie}, DocILE \cite{simsa2023docile}) and an evaluation corpus of 50 diverse Indian GST tax invoices encompassing multi-page PDFs, camera scans, mixed intra/inter-state supplies, multi-rate taxes (0\%--28\%), Section 17(5) blocked items, and borderless tables.

\begin{figure}[htbp]
\centerline{\includegraphics[width=0.88\columnwidth]{../../report_figures/fig5_ablation_benchmarks.png}}
\caption{Comparative performance across experimental ablation configurations (E1 to E6).}
\label{fig:ablation}
\end{figure}

\begin{table}[htbp]
\caption{Experimental Ablation Results (Configurations E1 to E6)}
\begin{center}
\begin{tabularx}{\columnwidth}{lXX}
\toprule
\textbf{Config} & \textbf{Architectural Subsystems} & \textbf{Table TEDS / Balance Pass} \\
\midrule
E1 (Baseline) & Raw SmolVLM-Instruct & 68.4\% / 0.0\% \\
E2 (Heuristic) & SmolVLM + Regex Extraction & 79.2\% / 12.5\% \\
E3 (Table Corr) & Stage 3 Table Corrector & 91.5\% / 46.8\% \\
E4 (Integrated) & End-to-End Extraction Pipeline & 95.0\% / 78.4\% \\
E5 (Accounting) & COA + AS 10 + GST POS & 97.8\% / 94.2\% \\
\textbf{E6 (Complete)} & \textbf{Full System + Invariant Suite} & \textbf{98.4\% / 100.0\%} \\
\bottomrule
\end{tabularx}
\label{tab:ablation}
\end{center}
\end{table}

\subsection{Ablation Analysis and Discussion}
Figure \ref{fig:ablation} and Table \ref{tab:ablation} detail results across all six configurations. Baseline E1 suffered severe table errors (Table TEDS 68.4\%, invariant pass rate 0.0\%) due to footer rows parsed as items. Heuristics (E2) raised invoice F1 to 79.2\% but missed layout shifts. The Table Layout Corrector (E3) produced the single largest structural gain, boosting Table TEDS to 91.5\% and reducing footer contamination by 94.2\%. Accounting Engine (E5) achieved 98.2\% capitalization accuracy and 97.4\% GST routing precision. Full production configuration E6 achieved 98.4\% Table TEDS, 97.8\% overall field F1, and a 100.0\% invariant pass rate with zero balance violations.

\begin{figure}[htbp]
\centerline{\includegraphics[width=0.88\columnwidth]{../../report_figures/fig6_test_suite_coverage.png}}
\caption{Automated pytest execution distribution across 35 unit and invariant test cases.}
\label{fig:test_suite}
\end{figure}

\begin{table}[htbp]
\caption{Pipeline Latency and Subsystem Profiling ($N = 100$ Runs)}
\begin{center}
\begin{tabularx}{\columnwidth}{lXX}
\toprule
\textbf{Stage} & \textbf{Subsystem Operation} & \textbf{Mean Latency (ms)} \\
\midrule
Stage 1 & PDF Ingestion \& Preprocessing & 148.2 ms \\
Stage 2 & Vector Token Normalization & 62.4 ms \\
Stage 3 & Table Layout Correction (Yao et al.) & 214.6 ms \\
Stage 4 & Line Item / Footer Decoupling & 86.1 ms \\
Stage 5 & Accounting Interpretation (AS 10/GST) & 184.5 ms \\
Stage 6 & Invariant Validation (6 Invariants) & 42.8 ms \\
Stage 7 & Authoritative Ledger Commitment & 104.0 ms \\
\midrule
\textbf{Total} & \textbf{End-to-End Autonomous Pipeline} & \textbf{842.6 ms} \\
\bottomrule
\end{tabularx}
\label{tab:latency}
\end{center}
\end{table}

\subsection{Test Suite Coverage \& Latency Profiling}
The system was evaluated against 35 automated pytest cases (Fig. \ref{fig:test_suite}) spanning Table Extraction, Accounting Rules, GST POS, Blocked Credits, Invariants, and Ledger Postings. All 35 tests passed in 7.04s (100\% pass rate). Table \ref{tab:latency} details latency across 100 runs on an Intel i7-13700H (32GB RAM); mean pipeline latency is 842.6 ms (p95 = 1120.4 ms), delivering a $>$100x speedup over manual entry.

\section{Case Study \& Numerical Walkthrough}
We trace a real B2B invoice (INV-2026-DEL-089) issued by ``Techtronics Solutions Pvt Ltd'' (GSTIN: 07AABCT1234D1Z2, Delhi) to a client in Noida, UP (POS Code: 09). The invoice comprises three heterogeneous items:
\begin{itemize}
    \item Line 1: ``Dell XPS 15 Laptop'' --- Qty: 1, Price: INR 1,15,000.00.
    \item Line 2: ``Logitech MX Master Mouse'' --- Qty: 2, Price: INR 1,250.00.
    \item Line 3: ``Team Lunch Catering'' --- Qty: 1, Price: INR 3,500.00.
    \item Tax: 18.0\% GST. Subtotal: INR 1,21,000; Tax: INR 21,780; Total: INR 1,42,780.
\end{itemize}

\subsection{Pipeline Execution and Cognitive Interpretation}
Line items are extracted and summaries trapped. Line 1 exceeds $\tau_{\text{cap}}$ (INR 10,000) and is capitalized as Non-Current Asset (Account 1060). Line 2 is $\le \tau_{\text{cap}}$ and is expensed (Account 5010). Line 3 is classified as Staff Welfare Expense (Account 5060). Because Supplier State (07) $\neq$ POS State (09), the transaction is inter-state (IGST). For Lines 1--2, tax is claimable ITC (Account 1083, INR 21,150.00). For Line 3, catering is blocked under Section 17(5); tax (INR 630.00) is routed to Ineligible GST Expense (Account 5080).

\begin{table}[htbp]
\caption{Balanced Journal Entry Posting for Case Study Invoice}
\begin{center}
\begin{tabularx}{\columnwidth}{llX}
\toprule
\textbf{Code} & \textbf{Account Title \& Classification} & \textbf{Debit / Credit (INR)} \\
\midrule
1060 & Computer Equipment (PPE) & Dr. INR 1,15,000.00 \\
5010 & Office Supplies Expense & Dr. INR 2,500.00 \\
5060 & Staff Welfare Expense & Dr. INR 3,500.00 \\
1083 & Input IGST (Tax Credit Asset) & Dr. INR 21,150.00 \\
5080 & Ineligible GST Expense (Sec 17(5)) & Dr. INR 630.00 \\
2010 & Accounts Payable (Techtronics) & Cr. INR 1,42,780.00 \\
\midrule
\textbf{Total} & \textbf{Double-Entry Exact Balance ($\Delta_{\text{bal}} = 0.00$)} & \textbf{INR 1,42,780.00} \\
\bottomrule
\end{tabularx}
\label{tab:journal}
\end{center}
\end{table}

Stage 6 verifies all invariants and confirms double-entry balance (Table \ref{tab:journal}):
$\text{Debits} = 1,15,000 + 2,500 + 3,500 + 21,150 + 630 = \text{INR } 1,42,780.00$;
$\text{Credits} = \text{INR } 1,42,780.00$;
$\Delta_{\text{bal}} = 0.000000$. Committed to ledger with zero floating-point drift.

\section{Discussion and Limitations}
A core design decision is placing deterministic invariant verification downstream of multimodal AI extraction. Deep learning models predicting ledger postings risk hallucinations and lack auditability. By decoupling extraction from deterministic accounting logic, probabilistic front-end extraction is paired with a zero-tolerance provable ledger back-end.

Current limitations include: 1) Severe physical scan degradation ($<$150 DPI) lowering OCR confidence; 2) Highly irregular nested or multi-page tabular geometries; 3) Lexicon maintenance for Section 17(5) blocked terms; and 4) Current prototype scope focusing on Accounts Payable (AP) vendor invoices rather than Accounts Receivable or automated three-way purchase order matching.

\section{Conclusion and Future Work}
We presented an end-to-end Multimodal Document Intelligence and Double-Entry Accounting Framework for autonomous accounting book generation. By bridging the semantic chasm between document extraction and bookkeeping, the framework unifies template-free extraction, table layout correction (Yao et al., 2026), AS 10 capitalization, statutory GST routing, and a 6-point invariant suite. The system achieves 98.4\% Table TEDS, 100\% exact double-entry balance, and passes 35/35 automated tests in 7.04s.

Future work will explore on-premise quantized VLM deployment (e.g., Qwen2.5-VL), automated three-way matching (PO, GRN, invoice), bi-directional Tally Prime and Zoho Books XML connectors, and multi-currency foreign exchange revaluation.

\bibliographystyle{IEEEtran}
\bibliography{references}

\end{document}
